\documentclass[11pt]{article}
\usepackage[spanish,es-noshorthands]{babel}
\usepackage[utf8]{inputenc}
\usepackage[T1]{fontenc}
\usepackage{amsmath,amssymb,amsthm,mathtools,booktabs,longtable,array,geometry,enumitem,xcolor}
\usepackage[hidelinks]{hyperref}
\geometry{margin=2.4cm}
\numberwithin{equation}{section}
\newtheorem{proposition}{Proposición}
\theoremstyle{remark}
\newtheorem*{remark}{Observación}

% --- notación vectorial y de operadores -------------------------------------
\newcommand{\bu}{\mathbf{u}}
\newcommand{\bK}{\mathbf{K}}
\newcommand{\bD}{\mathbf{D}}
\newcommand{\bn}{\mathbf{n}}
\newcommand{\bb}{\mathbf{b}}
\newcommand{\bd}{\mathbf{d}}
\newcommand{\bx}{\mathbf{x}}
\newcommand{\geps}{g_\varepsilon}
\newcommand{\im}{\mathrm{im}}
\newcommand{\vr}{\varrho}              % tasa de almacenamiento S_s \partial_t h
\newcommand{\pA}{\psi^{\mathrm{A}}}    % adjunto ADR
\newcommand{\pM}{\psi^{\mathrm{M}}}    % adjunto MRMT (fase móvil)
\newcommand{\pim}{\psi_{\im,r}}          % adjunto inmóvil
\newcommand{\ph}[1]{\psi_h^{\mathrm{#1}}}  % adjunto del flujo
\newcommand{\dd}{\,\mathrm{d}}
\newcommand{\Gin}{\Gamma_{\mathrm{in}}}
\newcommand{\Gout}{\Gamma_{\mathrm{out}}}
\newcommand{\Gwall}{\Gamma_{\mathrm{wall}}}

\title{Transporte y absorci\'on de contaminantes en un acu\'ifero:\\
modelos ADR y MRMT, formulaci\'on variacional, funcionales de Lagrange\\
y sistemas adjuntos}
\author{Tesis doctoral --- control \'optimo de la contaminaci\'on en acu\'iferos profundos}
\date{31 de octubre de 2026}

\begin{document}
\maketitle

\begin{abstract}
Se formulan los modelos de advecci\'on--dispersi\'on--reacci\'on (ADR) y de
transferencia de masa multitasa (MRMT) para el soluto en un acu\'ifero con un
pozo de extracci\'on. Se parte del balance conservativo de soluto y de la
ecuaci\'on de flujo; la divergencia de la velocidad de Darcy se despeja del
flujo, de modo que la extracci\'on del pozo queda incorporada de forma impl\'icita
en la ecuaci\'on de transporte, que conserva el t\'ermino de almacenamiento
$-S_s\,\partial_th\,C$. Se verifican las condiciones de frontera del
Cap\'itulo~2 de la tesis, se deducen las formulaciones d\'ebiles, se construye
un funcional de Lagrange independiente para cada modelo y se derivan paso a paso
los sistemas adjuntos (transporte, regiones inm\'oviles y flujo), sus condiciones
de frontera y terminales, y el gradiente respecto al caudal $Q$ y a la
profundidad del pozo $z_p$. Se describe la correspondencia con las
implementaciones RBF-FD (BFR) y de elementos finitos (FEniCSx) y se detalla qu\'e
verificaciones se han realizado.
\end{abstract}

\tableofcontents

%===========================================================
\section{Introducci\'on y descripci\'on del problema f\'isico}
%===========================================================

Un acu\'ifero confinado se representa por un corte vertical $\Omega\subset\mathbb{R}^2$
(las magnitudes de flujo son por unidad de espesor perpendicular al plano). El agua
subterr\'anea entra por la cara derecha, sale por la izquierda y se extrae mediante un
pozo situado en $(x_p,z_p)$ con caudal $Q(t)\ge0$. Un contaminante, presente como
condici\'on inicial y/o fuente $W_+$, es transportado por la velocidad de Darcy,
se dispersa, y en el modelo MRMT se intercambia con $N_r$ regiones inm\'oviles.
El objetivo del control \'optimo es elegir $(Q,z_p)$ que minimicen la
concentraci\'on de contaminante captada por el pozo y el costo econ\'omico.

El documento deduce el modelo directo, su formulaci\'on d\'ebil, los funcionales de Lagrange y los
sistemas adjuntos de los dos modelos, ADR y MRMT, de forma independiente.

%===========================================================
\section{Hip\'otesis del modelo}
%===========================================================

\begin{enumerate}[label=(H\arabic*)]
\item \textbf{Flujo.} Acu\'ifero confinado, ley de Darcy $\bu=-\bK\nabla h$ con
  $\bK=\mathrm{diag}(K_x,K_z)$ definido positivo y acotado, $S_s\ge\sigma_0>0$.
\item \textbf{Porosidad y retardo.} $\phi=\phi(\bx)$ y $R$ no dependen del tiempo; toda
  la variaci\'on temporal del almacenamiento de agua se expresa con $S_s\,\partial_th$.
\item \textbf{Pozo.} El pozo es un sumidero volum\'etrico de agua,
  $f=-Q(t)\,\geps(\bx;z_p)$, con $\geps\ge0$, $\int_\Omega\geps\dd\bx=1$.
\item \textbf{Soluto.} El agua extra\'ida lleva el soluto con la concentraci\'on local:
  retira masa a raz\'on $fC$ por unidad de volumen.
\item \textbf{Dispersi\'on.} $\bD=\bD(\bu)=\mathrm{diag}(D_x,D_z)$ definida positiva
  (Cap.~2, ec.~(2.4)), $D_d>0$.
\item \textbf{MRMT.} Las regiones inm\'oviles no se transportan; el intercambio es lineal,
  $\lambda_r(C_m-C_{\im,r})$, y $\phi_{\im,r}$ es constante en el tiempo.
\item \textbf{Regularidad.} $\bu\in L^\infty(Q_T)^2$, $\vr=S_s\partial_th\in L^\infty(Q_T)$,
  $\nabla\!\cdot\bu\in L^2(Q_T)$.
\end{enumerate}

%===========================================================
\section{Definici\'on del dominio y notaci\'on}
%===========================================================

$\Omega\subset\mathbb{R}^2$ es abierto, acotado y de frontera Lipschitz;
$Q_T=\Omega\times(0,T)$. La frontera se descompone en
$\partial\Omega=\overline{\Gin}\cup\overline{\Gout}\cup\overline{\Gwall}$:
$\Gin$ es la cara derecha (entrada), $\Gout$ la izquierda (salida) y $\Gwall$ las caras superior e inferior
(impermeables); $\bn$ es la normal unitaria exterior. $\langle\cdot,\cdot\rangle$ es el
producto de dualidad y $(\cdot,\cdot)$ el de $L^2(\Omega)$.

\paragraph{Diccionario de notaci\'on.} Cada letra designa una sola magnitud. Los
modelos se distinguen por el sub/superíndice: $C$ es la concentraci\'on del modelo ADR;
$C_m$ y $C_{\im,r}$ son las del MRMT; $\pA,\ph{A}$ son los adjuntos del ADR y
$\pM,\pim,\ph{M}$ los del MRMT. El tiempo de integraci\'on auxiliar se denota $\tau$.

{\small
\begin{longtable}{@{}p{1.9cm}p{4.6cm}p{4.2cm}p{1.9cm}p{1.2cm}@{}}
\toprule
\textbf{S\'imbolo} & \textbf{Definici\'on} & \textbf{Significado f\'isico} & \textbf{Unidades} & \textbf{Modelo}\\
\midrule
\endhead
$h(\bx,t)$ & $h\in\mathcal{W}(0,T;V,V')$ & carga hidr\'aulica (estado) & m & ambos\\
$\bu$ & $-\bK\nabla h$ & velocidad de Darcy & m/s & ambos\\
$\bK$ & $\mathrm{diag}(K_x,K_z)$ & conductividad hidr\'aulica & m/s & ambos\\
$S_s$ & constante $>0$ & almacenamiento espec\'ifico & m$^{-1}$ & ambos\\
$\vr$ & $S_s\,\partial_th$ & tasa de almacenamiento de agua & s$^{-1}$ & ambos\\
$f$ & $-Q\,\geps$ & extracci\'on volum\'etrica (agua) & s$^{-1}$ & ambos\\
$f_{\rm in}$ & dato, $\ge0$ & flujo entrante en $\Gin$ & m/s & ambos\\
$Q(t)$ & control, $L^2(0,T)$ & caudal por unidad de espesor & m$^2$/s & ambos\\
$z_p$ & control escalar & profundidad del pozo & m & ambos\\
$\geps$ & kernel gaussiano, $\int\geps=1$ & delta regularizada del pozo & m$^{-2}$ & ambos\\
$C$ & $C\in\mathcal{W}(0,T;V_C,V_C')$ & concentraci\'on (estado ADR) & $[C]$ & ADR\\
$C_m$ & idem & concentraci\'on m\'ovil (estado MRMT) & $[C]$ & MRMT\\
$C_{\im,r}$ & $\in H^1(0,T;L^2)$ & concentraci\'on inm\'ovil, regi\'on $r$ & $[C]$ & MRMT\\
$\phi,\ R$ & dados & porosidad m\'ovil, factor de retardo & ---, --- & ambos\\
$\phi_{\im,r}$ & $\eta_r\phi$ & porosidad de la regi\'on $r$ & --- & MRMT\\
$\lambda_d$ & dato & decaimiento de primer orden & s$^{-1}$ & ADR\\
$\lambda_r$ & $\phi D_{{\rm eff},r}/L^2$ & coeficiente de transferencia & s$^{-1}$ & MRMT\\
$\bD(\bu)$ & ec.~\eqref{eq:dispersion} & dispersi\'on hidrodin\'amica & m$^2$/s & ambos\\
$W_+$ & dato, $\ge0$ & fuente de soluto & $[C]$/s & ambos\\
$C_0,\,C_{\im,r,0},\,h_0$ & datos iniciales en $L^2(\Omega)$ & estados en $t=0$ & ---& ---\\
$\pA,\ \pM$ & adjuntos de $C,\,C_m$ & multiplicadores de transporte & $[C]\,$s\,m$^{-2}$ & ADR, MRMT\\
$\pim$ & adjunto de $C_{\im,r}$ & multiplicador inm\'ovil & $[C]\,$s\,m$^{-2}$ & MRMT\\
$\ph{A},\ph{M}$ & adjuntos de $h$ & multiplicadores del flujo & $[C]^2$s\,m$^{-2}$ & ADR, MRMT\\
$J_m$ & $\int_0^T\!\!\int_\Omega\geps C^2$ (ADR; $C_m$ en MRMT) & calidad del agua captada & $[C]^2$s & ambos\\
$J_{\rm inst}$ & $\kappa_p(z_p-z_0)^2$ & costo de instalaci\'on & --- & ambos\\
$J_{\rm op}$ & $d_p\int_0^TQ^2\dd t$,\ $d_p=z_0-z_p$ & costo energ\'etico & --- & ambos\\
$\gamma,\alpha,\beta$ & pesos $>0$ & ponderaci\'on de $J_m,J_{\rm inst},J_{\rm op}$ & --- & ambos\\
$\bb,\bd$ & ec.~\eqref{eq:bb} & acoplamiento transporte$\to$flujo & ver texto & ambos\\
$\tau$ & variable de integraci\'on & tiempo auxiliar & s & ---\\
\bottomrule
\end{longtable}}

$[C]$ denota las unidades de concentraci\'on. Los dominios son bidimensionales y todas las
magnitudes de flujo son por unidad de espesor.

%===========================================================
\section{Ecuaciones de gobierno y simplificaci\'on del t\'ermino de extracci\'on}
\label{sec:gobierno}
%===========================================================

\paragraph{Flujo.}
\begin{equation}
S_s\,\partial_th-\nabla\!\cdot(\bK\nabla h)=f=-Q\,\geps,\qquad \bu=-\bK\nabla h.
\label{eq:flujo}
\end{equation}
Como $\nabla\!\cdot\bu=-\nabla\!\cdot(\bK\nabla h)$, la divergencia se despeja del flujo:
\begin{equation}
\nabla\!\cdot\bu=f-S_s\,\partial_th=-Q\,\geps-\vr .
\label{eq:div_u}
\end{equation}

\paragraph{Balance conservativo de soluto.} La masa por unidad de volumen es $\phi RC$, el flujo
es $\bu C-\bD\nabla C$ y, por (H4), el pozo retira masa a raz\'on $fC$ ($f<0$ en el pozo):
\begin{equation}
\phi R\,\partial_tC+\nabla\!\cdot(\bu\,C-\bD\nabla C)+\lambda_d\phi R\,C=W_++f\,C .
\label{eq:conservativa}
\end{equation}
(La forma MRMT se trata en la Secci\'on~\ref{sec:mrmt}; el procedimiento es el mismo.)

\paragraph{Paso 1: desarrollo sin descartar $\nabla\!\cdot\bu$.}
$\nabla\!\cdot(\bu C)=\bu\cdot\nabla C+C\,\nabla\!\cdot\bu$. Luego
\begin{equation}
\phi R\,\partial_tC+\bu\cdot\nabla C-\nabla\!\cdot(\bD\nabla C)+\lambda_d\phi RC+C\,\nabla\!\cdot\bu=W_++f\,C.
\label{eq:desarrollada}
\end{equation}

\paragraph{Paso 2: sustituci\'on de la divergencia.} Por \eqref{eq:div_u}, $C\nabla\!\cdot\bu=fC-\vr\,C$:
\begin{equation}
\phi R\,\partial_tC+\bu\cdot\nabla C-\nabla\!\cdot(\bD\nabla C)+\lambda_d\phi RC+fC-\vr\,C=W_++fC .
\end{equation}

\paragraph{Paso 3: cancelaci\'on.} El t\'ermino $fC$ aparece en ambos miembros y se cancela;
el t\'ermino de almacenamiento $\vr\,C$ no:
\begin{equation}
\boxed{\ \phi R\,\partial_tC+\bu\cdot\nabla C-\nabla\!\cdot(\bD\nabla C)+\lambda_d\phi R\,C-S_s\,\partial_th\;C=W_+\ }
\label{eq:N_ADR}
\end{equation}

\begin{proposition}[Equivalencia]\label{prop:equiv}
Sea $h$ soluci\'on de \eqref{eq:flujo} con (H7). Entonces $C$ resuelve \eqref{eq:conservativa} si y s\'olo si
resuelve \eqref{eq:N_ADR}, en el sentido de $L^2(Q_T)$ y en el d\'ebil.
\end{proposition}
\begin{proof}
Los pasos 1--3 son identidades algebraicas reversibles; del flujo s\'olo se usa \eqref{eq:div_u}.
Con (H7) y $C\in L^2(0,T;H^1)$, el producto $C\,\nabla\!\cdot\bu$ es integrable, de modo que se puede probar con $v\in H^1\cap L^\infty$ y extender por densidad.
En sentido d\'ebil, $\int(\nabla\!\cdot(\bu C))v=\oint(\bu\cdot\bn)Cv-\int C\,\bu\cdot\nabla v$ y la identidad \eqref{eq:div_u} se usa en el mismo orden.
\end{proof}

\paragraph{La extracci\'on est\'a representada de forma impl\'icita.} En \eqref{eq:N_ADR} no hay un t\'ermino
$\chi\,\geps C$, pero el proceso f\'isico subsiste: cerca del pozo $\nabla\!\cdot\bu=-Q\geps-\vr<0$ (el campo de Darcy
converge) y el t\'ermino $\bu\cdot\nabla C$ conduce el soluto hacia all\'i. Lo muestra el balance de masa.
Sea $m=\int_\Omega\phi RC$. Integrando \eqref{eq:N_ADR}, con $\int_\Omega\bu\cdot\nabla C=\oint(\bu\cdot\bn)C-\int_\Omega C\nabla\!\cdot\bu$,
\eqref{eq:div_u} y la condici\'on $(\bD\nabla C)\cdot\bn=0$ (Secci\'on~\ref{sec:bc}):
\begin{equation}
\frac{\dd m}{\dd t}=-\oint_{\partial\Omega}(\bu\cdot\bn)\,C+\int_\Omega W_+-\lambda_dm\;\underbrace{-\,Q\!\int_\Omega\geps C}_{\text{extracci\'on}}\;
\underbrace{-\,S_s\!\int_\Omega\partial_th\,C+S_s\!\int_\Omega\partial_th\,C}_{=\,0}.
\label{eq:balance}
\end{equation}
El primer sumando verde de la divergencia ($+\vr C$ al despejar) y el t\'ermino $-\vr C$ de \eqref{eq:N_ADR} se anulan; queda el
sumidero $-Q\int\geps C$ (el mismo $fC$ de \eqref{eq:conservativa}) \emph{sin} t\'ermino expl\'icito en la ecuaci\'on.

\begin{remark}[Unidades]
$\geps$ tiene unidades m$^{-2}$ y $f=-Q\geps$ las de $S_s\partial_th$ (s$^{-1}$); por tanto $Q$ es caudal por unidad de espesor (m$^2$/s). Todos los t\'erminos
de \eqref{eq:N_ADR} tienen unidades $[C]/$s: $\phi R\partial_tC$, $\bu\cdot\nabla C$ (m/s\,$\cdot[C]$/m), $\nabla\!\cdot(\bD\nabla C)$ (m$^2$/s\,$\cdot[C]$/m$^2$),
$\lambda_d\phi RC$, $\vr C$ (s$^{-1}[C]$).
\end{remark}

\begin{remark}[Prueba local]
Donde $C$ es localmente constante y $W_+=\lambda_d=0$, \eqref{eq:N_ADR} da $\phi R\,\partial_tC=\vr\,C$: sin almacenamiento
($\partial_th=0$) extraer agua de concentraci\'on uniforme no modifica $C$.
\end{remark}

%===========================================================
\section{Modelo directo ADR}
\label{sec:adr}
%===========================================================

\paragraph{Problema.} Hallar $(h,C)$ tales que, en $Q_T$,
\begin{align}
&S_s\,\partial_th-\nabla\!\cdot(\bK\nabla h)=-Q\,\geps, && \bu=-\bK\nabla h,\label{eq:adr_h}\\
&\phi R\,\partial_tC+\bu\cdot\nabla C-\nabla\!\cdot(\bD(\bu)\nabla C)+\lambda_d\phi R\,C-S_s\,\partial_th\,C=W_+,\label{eq:adr_C}
\end{align}
con las condiciones iniciales y de frontera de la Secci\'on~\ref{sec:bc}. Las incógnitas son $h$ y $C$; los parámetros físicos
$\bK,S_s,\phi,R,\alpha_L,\alpha_T,D_d,\lambda_d$; los controles $Q$ y $z_p$.

\paragraph{Forma conservativa y forma expandida.} La forma \eqref{eq:conservativa}, con la divergencia de $\bu$ y el
t\'ermino $fC$ a la derecha, equivale a \eqref{eq:adr_C} (Proposici\'on~\ref{prop:equiv}). Equivalencia requiere \eqref{eq:div_u},
que se cumple exactamente a nivel continuo. No se pierde ning\'un t\'ermino: $C\nabla\!\cdot\bu$ se sustituye por $fC-\vr C$, no se descarta.
Si se supusiera $\nabla\!\cdot\bu=0$ (flujo solenoidal, falso con pozo y almacenamiento) se perder\'ian tanto $fC$ como $\vr C$.

\paragraph{Tensor de dispersi\'on.}
\begin{equation}
D_x=D_d+\frac{\alpha_Lu_x^2+\alpha_Tu_z^2}{|\bu|},\quad D_z=D_d+\frac{\alpha_Tu_x^2+\alpha_Lu_z^2}{|\bu|}.
\label{eq:dispersion}
\end{equation}

%===========================================================
\section{Modelo directo MRMT}
\label{sec:mrmt}
%===========================================================

El balance conservativo de la fase m\'ovil y de cada regi\'on inm\'ovil es
\begin{equation}
\phi R\,\partial_tC_m+\nabla\!\cdot(\bu C_m-\bD\nabla C_m)+\sum_{r=1}^{N_r}\lambda_r(C_m-C_{\im,r})=W_++f\,C_m,\qquad
\phi_{\im,r}R\,\partial_tC_{\im,r}=\lambda_r(C_m-C_{\im,r}).
\end{equation}
El intercambio no contiene $\bu$ ni $f$, de modo que los Pasos 1--3 de la Secci\'on~\ref{sec:gobierno} se repiten literalmente
sobre la ecuaci\'on de $C_m$ ($fC_m$ se cancela y $\vr C_m$ permanece):
\begin{equation}
\boxed{\begin{aligned}
&\phi R\,\partial_tC_m+\bu\cdot\nabla C_m-\nabla\!\cdot(\bD\nabla C_m)+\sum_{r=1}^{N_r}\lambda_r\,(C_m-C_{\im,r})-S_s\,\partial_th\,C_m=W_+,\\
&\phi_{\im,r}R\,\partial_tC_{\im,r}=\lambda_r\,(C_m-C_{\im,r}),\qquad r=1,\dots,N_r.
\end{aligned}}
\label{eq:N_MRMT}
\end{equation}
El flujo \eqref{eq:adr_h} es el mismo.

\paragraph{Relaci\'on con ADR.} Con $N_r=0$ y a\~nadiendo el decaimiento $\lambda_d\phi RC_m$ se recupera \eqref{eq:adr_C}. El MRMT no
lleva decaimiento: su t\'ermino de reacci\'on es el intercambio con las regiones inm\'oviles.

\paragraph{Conservaci\'on de masa.} Sea $M=\int_\Omega\bigl(\phi RC_m+\sum_r\phi_{\im,r}RC_{\im,r}\bigr)$. Sumando las ecuaciones, el
intercambio se cancela exactamente (aparece $+\lambda_r(C_m-C_{\im,r})$ en la m\'ovil y $-\lambda_r(C_m-C_{\im,r})$ en la inm\'ovil) y, por el
mismo c\'alculo de \eqref{eq:balance},
\[
\frac{\dd M}{\dd t}=-\oint_{\partial\Omega}(\bu\cdot\bn)C_m+\int_\Omega W_+-Q\!\int_\Omega\geps\,C_m .
\]
La extracci\'on es la de \eqref{eq:balance}, con la concentraci\'on de la fase m\'ovil, y no hay p\'erdida de masa en el intercambio.

%===========================================================
\section{Condiciones iniciales y de frontera}
\label{sec:bc}
%===========================================================

\paragraph{Iniciales.} $h(\cdot,0)=h_0$, $C(\cdot,0)=C_0$ (ADR); $C_m(\cdot,0)=C_{m,0}$, $C_{\im,r}(\cdot,0)=C_{\im,r,0}$ (MRMT), con datos en $L^2(\Omega)$.

\paragraph{Flujo hidr\'aulico (Cap.~2, Sec.~2.2.6).}
\begin{align}
&(\bK\nabla h)\cdot\bn=0 &&\text{en }\Gwall\ (\text{caras superior e inferior, impermeables}),\label{eq:bc_wall}\\
&(\bK\nabla h)\cdot\bn=f_{\rm in}(z),\ f_{\rm in}\ge0 &&\text{en }\Gin\ (\text{cara derecha, entrada}),\label{eq:bc_in}\\
&\nabla(\bK\nabla h)\cdot\bn=0 &&\text{en }\Gout\ (\text{cara izquierda, salida libre}).\label{eq:bc_out}
\end{align}
La condici\'on \eqref{eq:bc_out} es de segundo orden: no fija la carga ni el flujo, y anula la variaci\'on normal del flujo de Darcy en la salida.
En \eqref{eq:bc_in} la normal exterior en $\Gin$ es $(+1,0)$ y $\bu\cdot\bn=-f_{\rm in}\le0$: el flujo entra.\footnote{El
Cap\'itulo~2 (versi\'on \texttt{ChaptersD}) escribe $-(\bK\nabla h)\cdot\bn=f_{\rm in}$; con $\bu=-\bK\nabla h$ y $f_{\rm in}\ge0$ eso dar\'ia
$\bu\cdot\bn\ge0$ (salida). El signo de \eqref{eq:bc_in} es el de la implementaci\'on (FEniCS: $-(\bK\nabla h)\cdot\bn=h_{\rm inlet}=-f_{\rm in}$;
BFR: fila $-K_x\partial_xh=h_{\rm inlet}$) y el de \texttt{ChaptersR}. Se registra en el informe de auditor\'ia.}
En $\Gwall$, $\bu\cdot\bn=0$.

\paragraph{Contaminante.} La frontera es libre en el sentido del Cap.~2 (ecuaci\'on de frontera en la discretizaci\'on RBF-FD): flujo
dispersivo normal nulo
\begin{equation}
(\bD\nabla C)\cdot\bn=0\quad\text{en }\partial\Omega=\Gin\cup\Gout\cup\Gwall,
\label{eq:bc_C}
\end{equation}
y de igual forma para $C_m$. Es una condici\'on sobre el flujo \emph{difusivo}: el flujo total $\bu C-\bD\nabla C$ no es nulo en $\Gin$ ni $\Gout$,
donde $\bu\cdot\bn\ne0$ y la advecci\'on cruza la frontera (por eso aparece $-\oint(\bu\cdot\bn)C$ en \eqref{eq:balance}); s\'olo en $\Gwall$,
donde $\bu\cdot\bn=0$, el flujo total es nulo. Las regiones inm\'oviles no requieren condiciones de frontera (EDO local).

\begin{remark}
El Cap\'itulo~2 (\texttt{ChaptersD}, ec.~(2.11)) enuncia en cambio $C_m=0$ en $\Gin\cup\Gout$ para el modelo continuo y declara que RBF-FD lo sustituye
por \eqref{eq:bc_C}; \texttt{ChaptersR} impone \eqref{eq:bc_C} en toda la frontera y las dos implementaciones (BFR y FEniCS) tambi\'en. Aqu\'i se adopta
\eqref{eq:bc_C}; la discrepancia se documenta en el informe de auditor\'ia.
\end{remark}

\begin{table}[h]
\centering\small
\begin{tabular}{p{2.2cm}p{1.5cm}p{4.3cm}p{4.2cm}p{2.5cm}}
\toprule
Frontera & Campo & Condici\'on & Interpretaci\'on & Referencia\\
\midrule
$\Gwall$ (sup., inf.) & $h$ & $(\bK\nabla h)\cdot\bn=0$ & impermeable, $\bu\cdot\bn=0$ & Cap.~2, ec.~(2.7)\\
$\Gin$ (derecha) & $h$ & $(\bK\nabla h)\cdot\bn=f_{\rm in}(z)$ & entrada prescrita, $f_{\rm in}\ge0$ & Cap.~2, ec.~(2.8) (signo, ver nota)\\
$\Gout$ (izquierda) & $h$ & $\nabla(\bK\nabla h)\cdot\bn=0$ & salida libre de 2.\textordmasculine\ orden & Cap.~2, ec.~(2.9)\\
$\Gwall$ & $C,C_m$ & $(\bD\nabla C)\cdot\bn=0$ & flujo dispersivo nulo (=total nulo) & Cap.~2, ec.~(2.10)\\
$\Gin,\Gout$ & $C,C_m$ & $(\bD\nabla C)\cdot\bn=0$ & dispersi\'on nula, advecci\'on libre & Cap.~2, Sec.~2.5.5 y \texttt{ChaptersR}\\
$\Omega$ & $C_{\im,r}$ & ninguna & EDO local & Cap.~2, ec.~(2.6)\\
\bottomrule
\end{tabular}
\end{table}

%===========================================================
\section{Formulaciones variacionales}
\label{sec:debil}
%===========================================================

\paragraph{Espacios.} $V=H^1(\Omega)$ para $h$; $V_C=H^1(\Omega)$ para $C,C_m$ (no hay Dirichlet); $L^2(\Omega)$ para $C_{\im,r}$.
Evoluci\'on: $\mathcal{W}(0,T;V,V')=\{v\in L^2(0,T;V):\partial_tv\in L^2(0,T;V')\}$. Regularidad adicional en $\Gout$ para la traza
$\bK\nabla h\cdot\bn$: $h\in H^2$ cerca de $\Gout$ (Cap.~2).

\paragraph{Flujo.} Se multiplica \eqref{eq:adr_h} por $v\in V$ y se integra. Green: $-\int\nabla\!\cdot(\bK\nabla h)v=\int\bK\nabla h\cdot\nabla v-\oint(\bK\nabla h\cdot\bn)v$.
El t\'ermino de frontera se descompone: en $\Gwall$ se anula por \eqref{eq:bc_wall}; en $\Gin$ vale $\int_{\Gin}f_{\rm in}v$ por \eqref{eq:bc_in}; en $\Gout$ no hay dato y se
conserva:
\begin{equation}
\int_\Omega S_s\partial_th\,v+\underbrace{\int_\Omega\bK\nabla h\cdot\nabla v-\int_{\Gout}(\bK\nabla h\cdot\bn)v}_{a_h(h,v)}
=\underbrace{-Q\!\int_\Omega\geps v+\int_{\Gin}f_{\rm in}v}_{\ell_h(v;Q,z_p)}\quad\forall v\in V .
\label{eq:debil_h}
\end{equation}

\paragraph{Transporte ADR.} Se multiplica \eqref{eq:adr_C} por $v\in V_C$:
$-\int\nabla\!\cdot(\bD\nabla C)v=\int\bD\nabla C\cdot\nabla v-\oint(\bD\nabla C\cdot\bn)v$. La integral de frontera se anula en las tres partes por \eqref{eq:bc_C}
(no porque $v=0$): es la condici\'on natural. Resulta
\begin{equation}
\begin{aligned}
&\int_\Omega\phi R\,\partial_tC\,v+a_C^{\mathrm A}(C,v;h)=\int_\Omega W_+v\quad\forall v\in V_C,\\
&a_C^{\mathrm A}(C,v;h)=\int_\Omega\bD(\bu)\nabla C\cdot\nabla v+\int_\Omega(\bu\cdot\nabla C)v+\lambda_d\!\int_\Omega\phi RCv-\int_\Omega\vr\,Cv .
\end{aligned}
\label{eq:debil_adr}
\end{equation}

\paragraph{Transporte MRMT.}
\begin{equation}
\begin{aligned}
&\int_\Omega\phi R\,\partial_tC_m\,v+a_C^{\mathrm M}(C_m,v;h)-\sum_r\lambda_r(C_{\im,r},v)=\int_\Omega W_+v\quad\forall v\in V_C,\\
&\int_\Omega\phi_{\im,r}R\,\partial_tC_{\im,r}\,w_r=\lambda_r(C_m-C_{\im,r},w_r)\quad\forall w_r\in L^2,\ r=1,\dots,N_r,
\end{aligned}
\label{eq:debil_mrmt}
\end{equation}
con $a_C^{\mathrm M}(C_m,v;h)=\int\bD\nabla C_m\cdot\nabla v+\int(\bu\cdot\nabla C_m)v+c_0\!\int C_mv-\int\vr C_mv$, $c_0=\sum_r\lambda_r$.
Las contribuciones de todas las ecuaciones de almacenamiento y los acoplamientos m\'ovil--inm\'ovil son los dos sumandos con $\lambda_r$.

\paragraph{Bien-planteamiento y energ\'ia.} Probando con $v=C$ en \eqref{eq:debil_adr}: $\int(\bu\cdot\nabla C)C=\tfrac12\oint(\bu\cdot\bn)C^2-\tfrac12\int(\nabla\!\cdot\bu)C^2$, y
con \eqref{eq:div_u} el coeficiente de $C^2$ en el volumen es $\tfrac12(Q\geps+\vr)-\vr=\tfrac12(Q\geps-\vr)$, acotado por (H7). El t\'ermino de frontera
$\tfrac12\oint(\bu\cdot\bn)C^2$ ya no se anula (no hay $C=0$): se acota con la desigualdad de traza, $\bigl|\oint(\bu\cdot\bn)C^2\bigr|\le\|\bu\|_\infty\|C\|^2_{L^2(\partial\Omega)}\le\epsilon\|\nabla C\|^2+c_\epsilon\|C\|^2$,
y se absorbe en la dispersi\'on porque $\bD\ge D_d>0$. Gronwall da la estimaci\'on de energ\'ia y existencia y unicidad por Galerkin. En MRMT el intercambio es monótono: $\sum_r\lambda_r(C_m-C_{\im,r})^2\ge0$. El
flujo es de Lions con la salvedad de $\Gout$ (Cap.~2).

%===========================================================
\section{Problema de optimizaci\'on y funcional objetivo}
\label{sec:opt}
%===========================================================

El problema (Cap.~3) es $\min_{(Q,z_p)\in\mathcal{U}_{ad}}\hat J(Q,z_p)=J(C(Q,z_p),Q,z_p)$ sujeto a las ecuaciones de estado, con
\[
J=\gamma\,J_m+\alpha\,J_{\rm inst}+\beta\,J_{\rm op},\quad
J_m=\int_0^T\!\!\int_\Omega\geps(\bx;z_p)\,C^2\dd\bx\dd t,\quad J_{\rm inst}=\kappa_p(z_p-z_0)^2,\quad J_{\rm op}=d_p\int_0^TQ^2\dd t,
\]
$d_p=z_0-z_p>0$, y $C$ se sustituye por $C_m$ en MRMT. $\mathcal{U}_{ad}=\{(Q,z_p):Q\in L^2(0,T),\,0\le Q\le Q_{\max},\,z_p\in[z_p^{\min},z_p^{\max}]\}$.
Las variables de estado son $(h,C)$ o $(h,C_m,C_{\im,r})$; las restricciones, las ecuaciones de la Secci\'on~\ref{sec:debil} con sus condiciones iniciales y de frontera.
No hay observaciones experimentales: $J_m$ es una medida de exposici\'on en el entorno del pozo con el kernel $\geps$, no un ajuste a datos.

%===========================================================
\section{Funcional de Lagrange del modelo ADR}
\label{sec:L_ADR}
%===========================================================

Convenci\'on: los residuos se suman a $J$ (multiplicadores con signo $+$). Las condiciones iniciales $h(0)=h_0$, $C(0)=C_0$ se imponen en el conjunto de estados admisibles
(sus variaciones cumplen $\delta h(0)=\delta C(0)=0$); las terminales de los multiplicadores surgen de integrar por partes en el tiempo.
Los multiplicadores son $\ph{A}\in L^2(0,T;V)\cap H^1(0,T;V')$ y $\pA\in L^2(0,T;V_C)\cap H^1(0,T;V_C')$.
\begin{multline}
\mathcal{L}^{\mathrm A}=J+\int_0^T\Bigl[\langle S_s\partial_th,\ph{A}\rangle+a_h(h,\ph{A})-\ell_h(\ph{A};Q,z_p)
+\langle\phi R\,\partial_tC,\pA\rangle+a_C^{\mathrm A}(C,\pA;h)-(W_+,\pA)\Bigr]\dd t .
\label{eq:L_ADR}
\end{multline}
\paragraph{Verificaci\'on.} (i) $\delta_{\ph{A}}\mathcal{L}^{\rm A}=0$ para todo $\ph{A}$ recupera \eqref{eq:debil_h}; $\delta_{\pA}\mathcal{L}^{\rm A}=0$ recupera \eqref{eq:debil_adr}.
(ii) Cada multiplicador est\'a asociado a su ecuaci\'on ($\ph{A}$ al flujo, $\pA$ al transporte). (iii) Los t\'erminos de frontera de $\Gin$, $\Gout$ y $\Gwall$ est\'an incluidos en $a_h$ y $\ell_h$
y los del transporte son nulos por \eqref{eq:bc_C}. (iv) $J$ est\'a definido sobre $C$ y los controles. (v) No aparece ning\'un t\'ermino expl\'icito de extracci\'on en la parte de transporte: el control entra en
$\ell_h$ y en $J$.

%===========================================================
\section{Derivaci\'on del adjunto ADR}
\label{sec:adj_ADR}
%===========================================================

\subsection{Variaci\'on respecto de $C$}
Sea $\eta\in L^2(0,T;V_C)\cap H^1(0,T;V_C')$, $\eta(0)=0$. El Lagrangiano depende de $C$ por $J_m$, por $\phi R\partial_tC$ y por $a_C^{\mathrm A}$ (a trav\'es de $C$, no de $h$).
\begin{enumerate}[label=(\alph*)]
\item $\gamma J_m$: $2\gamma\int_0^T(\geps C,\eta)\dd t$.
\item Tiempo: $\int_0^T\langle\phi R\,\partial_t\eta,\pA\rangle\dd t=\bigl[\phi R(\eta,\pA)\bigr]_0^T-\int_0^T\langle\phi R\,\partial_t\pA,\eta\rangle\dd t=-\int_0^T\langle\phi R\,\partial_t\pA,\eta\rangle\dd t$,
 con $\eta(0)=0$ y la condici\'on terminal $\pA(T)=0$ que se impone para anular el t\'ermino en $T$.
\item $a_C^{\mathrm A}(\eta,\pA;h)$.
\end{enumerate}
Anulando la suma para todo $\eta$ se obtiene la forma d\'ebil del adjunto: $\pA(t)\in V_C$, $\pA(T)=0$,
\begin{equation}
-\langle\phi R\,\partial_t\pA,\eta\rangle+a_C^{\mathrm A}(\eta,\pA;h)=-2\gamma(\geps C,\eta)\qquad\forall\eta\in V_C .
\label{eq:adj_debil_A}
\end{equation}

\subsection{Forma fuerte e integraci\'on por partes espacial}
Con $\eta\in H^1$ arbitraria (sin anularse en la frontera) y $\bD$ sim\'etrico:
\begin{align*}
\int_\Omega\bD\nabla\eta\cdot\nabla\pA&=-\int_\Omega\eta\,\nabla\!\cdot(\bD\nabla\pA)+\oint_{\partial\Omega}\eta\,(\bD\nabla\pA)\cdot\bn,\\
\int_\Omega(\bu\cdot\nabla\eta)\pA&=-\int_\Omega\eta\,\nabla\!\cdot(\bu\pA)+\oint_{\partial\Omega}\eta\,(\bu\cdot\bn)\pA .
\end{align*}
El operador advectivo se transpone con cambio de signo y aparece $\nabla\!\cdot(\bu\pA)$, no $\bu\cdot\nabla\pA$; el difusivo es autoadjunto; el t\'ermino temporal cambia de signo;
los t\'erminos de reacci\'on ($\lambda_d\phi R$ y $-\vr$) son multiplicativos y se conservan. Resulta
\begin{equation}
\boxed{\begin{aligned}
&-\phi R\,\partial_t\pA-\nabla\!\cdot(\bD\nabla\pA)-\nabla\!\cdot(\bu\,\pA)+\lambda_d\phi R\,\pA-S_s\partial_th\,\pA=-2\gamma\,\geps\,C\quad\text{en }Q_T,\\
&\pA(\cdot,T)=0,\qquad(\bD\nabla\pA+\bu\,\pA)\cdot\bn=0\ \ \text{en }\partial\Omega .
\end{aligned}}
\label{eq:adj_fuerte_A}
\end{equation}
La condici\'on de frontera es \emph{natural}: es lo que queda al anular el coeficiente de $\eta|_{\partial\Omega}$ ($\eta$ arbitraria en la traza). \emph{No} es la del problema directo: es de flujo total nulo (difusivo
m\'as advectivo) para $\pA$, y s\'olo coincide con $(\bD\nabla\pA)\cdot\bn=0$ en $\Gwall$ ($\bu\cdot\bn=0$). El operador adjunto no es id\'entico al directo.

\subsection{Forma no conservativa}
$-\nabla\!\cdot(\bu\pA)=-\bu\cdot\nabla\pA-\pA\,\nabla\!\cdot\bu$ y, por \eqref{eq:div_u}, $-\pA\nabla\!\cdot\bu=Q\geps\,\pA+\vr\,\pA$. Entonces
\begin{equation}
-\phi R\,\partial_t\pA-\nabla\!\cdot(\bD\nabla\pA)-\bu\cdot\nabla\pA+\lambda_d\phi R\,\pA+Q\,\geps\,\pA+\underbrace{\vr\,\pA-\vr\,\pA}_{=\,0}=-2\gamma\,\geps\,C .
\label{eq:adj_noconservativa_A}
\end{equation}
Los dos t\'erminos de almacenamiento se cancelan: el pozo aparece en el adjunto solo como $Q\geps\,\pA$, consecuencia de que $\nabla\!\cdot\bu\neq0$, no de un sumidero expl\'icito en el estado.
Ambas formas son equivalentes si y solo si \eqref{eq:div_u} (con regularidad (H7)) y la misma condici\'on de frontera.

\subsection{Variaci\'on respecto de $h$ y adjunto del flujo}
$h$ entra en $\mathcal{L}^{\mathrm A}$ por la ecuaci\'on de flujo, por $\bu=-\bK\nabla h$ (en $\bD(\bu)$ y en la advecci\'on) y por $\vr=S_s\partial_th$. Sea $\delta h\in V$, $\delta h(0)=0$, $\delta\bu=-\bK\nabla\delta h$.
\[
\partial_ha_C^{\mathrm A}[\delta h]=\int_\Omega(\delta\bu\cdot\nabla C)\pA+\int_\Omega\bigl(\partial_{\bu}\bD[\delta\bu]\nabla C\bigr)\!\cdot\!\nabla\pA-\int_\Omega S_s\partial_t\delta h\;C\,\pA
=-\int_\Omega\bK\bb\cdot\nabla\delta h-S_s\!\int_\Omega\partial_t\delta h\;C\,\pA,
\]
\begin{equation}
\bb=\pA\,\nabla C+\bd,\qquad\bd=\partial_xC\,\partial_x\pA\,\nabla_{\bu}D_x+\partial_zC\,\partial_z\pA\,\nabla_{\bu}D_z .
\label{eq:bb}
\end{equation}
Aqu\'i $\nabla_{\bu}D_x,\nabla_{\bu}D_z$ son los gradientes de \eqref{eq:dispersion} respecto de $\bu$. Integrando por partes en el tiempo con $\delta h(0)=0$ y $C\pA(T)=0$:
\[
-S_s\!\int_0^T\!\!\!\int_\Omega\partial_t\delta h\,C\pA=+S_s\!\int_0^T\!\!\!\int_\Omega\delta h\;\partial_t(C\pA),\qquad
\int_0^T\langle S_s\partial_t\delta h,\ph{A}\rangle=-\int_0^T\langle S_s\partial_t\ph{A},\delta h\rangle\ \ (\ph{A}(T)=0).
\]
Anulando $\partial_h\mathcal{L}^{\mathrm A}$ para toda $\delta h$: $\ph{A}(T)=0$ y
\begin{equation}
-\langle S_s\partial_t\ph{A},v\rangle+a_h(v,\ph{A})=\int_\Omega\bK\bb\cdot\nabla v-S_s\!\int_\Omega\partial_t(C\pA)\,v\qquad\forall v\in V .
\label{eq:adj_flujo_debil_A}
\end{equation}
Integrando por partes $a_h$ y el t\'ermino con $\bb$:
\begin{equation}
\boxed{-S_s\,\partial_t\ph{A}-\nabla\!\cdot(\bK\nabla\ph{A})=-\nabla\!\cdot(\bK\bb)-S_s\,\partial_t(C\,\pA)\ \ \text{en }Q_T,\qquad\ph{A}(T)=0.}
\label{eq:adj_flujo_A}
\end{equation}
El t\'ermino $-S_s\partial_t(C\pA)$ procede del almacenamiento en el transporte y acopla el adjunto del flujo con $C\pA$ incluso con $\bb=0$.

\subsection{Condiciones de frontera del adjunto del flujo}
El resto del Lagrangiano produce, con $\delta h$ y $\bK\nabla\delta h\cdot\bn$ como trazas independientes donde no hay dato,
\[
\oint\delta h\,\bK(\nabla\ph{A}-\bb)\cdot\bn-\int_{\Gout}(\bK\nabla\delta h\cdot\bn)\,\ph{A}=0 .
\]
En $\Gin$ el flujo es dato (Neumann): $\delta h$ libre $\Rightarrow\bK(\nabla\ph{A}-\bb)\cdot\bn=0$ (natural). En $\Gwall$ ocurre lo mismo. En $\Gout$ la forma $a_h$ contiene $-\int_{\Gout}(\bK\nabla h\cdot\bn)v$ y no es sim\'etrica:
$\delta h$ y $\bK\nabla\delta h\cdot\bn$ son independientes, y
\begin{equation}
\bK\nabla\ph{A}\cdot\bn=\bK\bb\cdot\bn\ \text{ en }\Gin\cup\Gwall,\qquad
\ph{A}=0\ \text{ y }\ \bK\nabla\ph{A}\cdot\bn=\bK\bb\cdot\bn\ \text{ en }\Gout .
\label{eq:bc_psih}
\end{equation}
En $\Gout$ el adjunto queda sobredeterminado (dos condiciones): no es un problema de contorno cl\'asico y la hip\'otesis de que el gradiente adjunto es la derivada del funcional reducido no se puede probar a nivel continuo
(Cap.~3, ``La frontera libre en el adjunto''). Por ello el adjunto se define a nivel discreto, donde la matriz de flujo y su transpuesta son invertibles, y se valida con el test de Taylor. Las condiciones \eqref{eq:bc_psih}
son consecuencia matem\'atica de la derivaci\'on y no hip\'otesis adicionales; el nulo en $\Gin\cup\Gwall$ es natural, el par en $\Gout$ procede del t\'ermino no sim\'etrico.

\subsection{Gradiente}
El control aparece en $\ell_h$ (a trav\'es de $Q$ y de $\geps(\cdot;z_p)$), en $J_m$ (kernel) y en $\beta J_{\rm op}$, $\alpha J_{\rm inst}$; no en el transporte. Con $\delta Q$:
\begin{equation}
\boxed{\ \nabla_Q\hat J(t)=2\beta\,d_p\,Q(t)+\int_\Omega\geps\,\ph{A}\dd\bx\ }
\label{eq:grad_Q}
\end{equation}
\begin{equation}
\boxed{\ \frac{\partial\hat J}{\partial z_p}=\gamma\!\int_0^T\!\!\!\int_\Omega\dot g_\varepsilon\,C^2+2\alpha\kappa_p(z_p-z_0)-\beta\!\int_0^TQ^2\dd t+\int_0^T\!\!Q(t)\!\int_\Omega\dot g_\varepsilon\,\ph{A}\dd\bx\dd t\ }
\label{eq:grad_z}
\end{equation}
con $\dot g_\varepsilon=\partial\geps/\partial z_p$ y $\partial d_p/\partial z_p=-1$. Con $Q$ constante a trozos, $\partial\hat J/\partial Q_n=\int_{I_n}\nabla_Q\hat J\dd t$.
\paragraph{Comprobaci\'on con las sensibilidades.} Para $\delta Q$: $S_s\partial_t\delta h-\nabla\!\cdot(\bK\nabla\delta h)=-\delta Q\,\geps$ y la linealizada del transporte tiene, adem\'as de $\delta\bu\cdot\nabla C$ y $\bD'[\delta\bu]$, el t\'ermino
$-S_s\partial_t\delta h\,C$ (sin fuente directa en $\delta Q$). Probando \eqref{eq:adj_debil_A} con $\delta C$ y la linealizada con $\pA$, se cancelan los t\'erminos de $\delta C$, y queda
$2\gamma\int_0^T(\geps C,\delta C)=-\iint\bK\bb\cdot\nabla\delta h+S_s\iint\delta h\,\partial_t(C\pA)$; probando \eqref{eq:adj_flujo_debil_A} con $v=\delta h$ y la ecuaci\'on de $\delta h$ con $\ph{A}$ se obtiene
$\int\delta Q\,(\int\geps\ph{A})\dd t$, el segundo t\'ermino de \eqref{eq:grad_Q}. El t\'ermino $-S_s\partial_t(C\pA)$ es el que cierra esta identidad.

%===========================================================
\section{Funcional de Lagrange del modelo MRMT}
\label{sec:L_MRMT}
%===========================================================

Multiplicadores: $\ph{M}$ (flujo), $\pM\in L^2(0,T;V_C)$ (fase m\'ovil), $\pim\in H^1(0,T;L^2(\Omega))$, $r=1..N_r$ (una por cada ecuaci\'on independiente).
\begin{multline}
\mathcal{L}^{\mathrm M}=J+\int_0^T\Bigl[\langle S_s\partial_th,\ph{M}\rangle+a_h(h,\ph{M})-\ell_h(\ph{M};Q,z_p)
+\langle\phi R\,\partial_tC_m,\pM\rangle+a_C^{\mathrm M}(C_m,\pM;h)\\
-\sum_{r=1}^{N_r}\lambda_r(C_{\im,r},\pM)-(W_+,\pM)
+\sum_{r=1}^{N_r}\bigl(\phi_{\im,r}R\,\partial_tC_{\im,r}+\lambda_r(C_{\im,r}-C_m),\ \pim\bigr)\Bigr]\dd t .
\label{eq:L_MRMT}
\end{multline}
Las variaciones respecto de los multiplicadores recuperan \eqref{eq:debil_h} y \eqref{eq:debil_mrmt}. El intercambio aparece con signo $+\lambda_r(C_m-C_{\im,r})$ en la ecuaci\'on m\'ovil (sumandos de $a_C^{\rm M}$ y $-\lambda_rC_{\im,r}$) y
$-\lambda_r(C_m-C_{\im,r})$ en la inm\'ovil, compatible con la conservaci\'on de masa de la Secci\'on~\ref{sec:mrmt}. $J$ depende \'unicamente de $C_m$.

%===========================================================
\section{Derivaci\'on del adjunto MRMT}
\label{sec:adj_MRMT}
%===========================================================

\subsection{Variaci\'on respecto de $C_m$}
Id\'entica a la del ADR con $c_0=\sum_r\lambda_r$ en lugar de $\lambda_d\phi R$ y con el acoplamiento: en $\mathcal{L}^{\rm M}$, $\delta C_m=\eta$ aporta (a) $2\gamma(\geps C_m,\eta)$,
(b) $-\langle\phi R\,\partial_t\pM,\eta\rangle$ (con $\eta(0)=0$, $\pM(T)=0$), (c) $a_C^{\rm M}(\eta,\pM;h)$, (d) de la ecuaci\'on inm\'ovil, $-\sum_r\lambda_r(\eta,\pim)$.
La derivaci\'on espacial es la de la Secci\'on~\ref{sec:adj_ADR}. Forma fuerte conservativa:
\begin{equation}
-\phi R\,\partial_t\pM-\nabla\!\cdot(\bD\nabla\pM)-\nabla\!\cdot(\bu\pM)+\sum_r\lambda_r(\pM-\pim)-S_s\partial_th\,\pM=-2\gamma\,\geps\,C_m,\qquad(\bD\nabla\pM+\bu\pM)\cdot\bn=0 .
\label{eq:adj_fuerte_M}
\end{equation}

\subsection{Variaci\'on respecto de $C_{\im,r}$}
$J$ no depende de $C_{\im,r}$. Con $\xi_r(0)=0$: el intercambio m\'ovil da $-\lambda_r(\xi_r,\pM)$ y la ecuaci\'on inm\'ovil da
$(\phi_{\im,r}R\,\partial_t\xi_r+\lambda_r\xi_r,\pim)$. Integrando por partes en el tiempo, $\int_0^T\phi_{\im,r}R(\partial_t\xi_r,\pim)=-\int_0^T\phi_{\im,r}R(\xi_r,\partial_t\pim)$ (con $\pim(T)=0$). La suma debe anularse para toda $\xi_r$:
\begin{equation}
\boxed{\ \phi_{\im,r}R\,\partial_t\pim=\lambda_r\,(\pim-\pM),\qquad\pim(\cdot,T)=0\ }
\label{eq:adj_inmovil}
\end{equation}
Resuelta hacia atr\'as con $\alpha_r=\lambda_r/(\phi_{\im,r}R)$: $\pim(t)=\alpha_r\int_t^Te^{-\alpha_r(\tau-t)}\pM(\tau)\dd\tau$, una memoria hacia el futuro. Ni el sumidero ni el almacenamiento actúan sobre las regiones inm\'oviles.
Como la aportaci\'on adjunta del acoplamiento m\'ovil es $\sum_r\lambda_r(\pM-\pim)$, su signo coincide con el del directo $\sum_r\lambda_r(C_m-C_{\im,r})$ (la transpuesta de un intercambio sim\'etrico).

\subsection{Forma no conservativa y sistema adjunto del transporte}
Igual que en ADR, $-\nabla\!\cdot(\bu\pM)=-\bu\cdot\nabla\pM+Q\geps\pM+\vr\pM$ y los verdes se cancelan:
\begin{equation}
\left\{\begin{aligned}
&-\phi R\,\partial_t\pM-\nabla\!\cdot(\bD\nabla\pM)-\bu\cdot\nabla\pM+Q\,\geps\,\pM+\sum_r\lambda_r(\pM-\pim)=-2\gamma\,\geps\,C_m,\\
&\phi_{\im,r}R\,\partial_t\pim=\lambda_r(\pim-\pM),\quad r=1..N_r,\\
&\pM(T)=\pim(T)=0,\qquad(\bD\nabla\pM+\bu\pM)\cdot\bn=0\ \text{ en }\partial\Omega .
\end{aligned}\right.
\label{eq:sistema_adjunto_M}
\end{equation}
La fuente depende \'unicamente de $C_m$ porque $J_m$ s\'olo depende de la fase m\'ovil; si $J$ dependiera de $C_{\im,r}$ aparecer\'ia una fuente en \eqref{eq:adj_inmovil}.

\subsection{Adjunto del flujo para MRMT}
$h$ entra igual (por $\bu$, $\bD(\bu)$ y $\vr$) con $C\to C_m$, $\pA\to\pM$. Entonces
\begin{equation}
-S_s\partial_t\ph{M}-\nabla\!\cdot(\bK\nabla\ph{M})=-\nabla\!\cdot(\bK\bb^{\mathrm M})-S_s\,\partial_t(C_m\pM),\quad\bb^{\mathrm M}=\pM\nabla C_m+\bd^{\mathrm M},\quad\ph{M}(T)=0,
\label{eq:adj_flujo_M}
\end{equation}
con las condiciones \eqref{eq:bc_psih} y el mismo gradiente \eqref{eq:grad_Q}--\eqref{eq:grad_z} reemplazando $(C,\ph{A})$ por $(C_m,\ph{M})$. Las regiones inm\'oviles no entran en el flujo.

%===========================================================
\section{Condiciones de frontera y terminales adjuntas}
\label{sec:bc_adj}
%===========================================================

\begin{table}[h]
\centering\small
\begin{tabular}{p{2.3cm}p{3.6cm}p{4.3cm}p{4.2cm}}
\toprule
Variable & Frontera & Condici\'on adjunta & Naturaleza / origen\\
\midrule
$\pA,\pM$ & $\partial\Omega$ & $(\bD\nabla\psi+\bu\psi)\cdot\bn=0$ & natural; resulta de $\oint\eta[(\bD\nabla\psi+\bu\psi)\cdot\bn]$ con $\eta$ libre\\
$\pim$ & --- & ninguna & EDO\\
$\ph{A},\ph{M}$ & $\Gin\cup\Gwall$ & $\bK\nabla\psi_h\cdot\bn=\bK\bb\cdot\bn$ & natural\\
$\ph{A},\ph{M}$ & $\Gout$ & $\psi_h=0$ y $\bK\nabla\psi_h\cdot\bn=\bK\bb\cdot\bn$ & del t\'ermino $-\int_{\Gout}(\bK\nabla h\cdot\bn)v$ de $a_h$; sobredeterminado\\
todas & $t=T$ & $\psi(T)=0$ & $J$ no tiene t\'ermino terminal (sólo integral en $(0,T)$)\\
\bottomrule
\end{tabular}
\end{table}
Condiciones esenciales: ninguna del transporte (los espacios son $H^1$). Condiciones homogéneas: terminales y las de $\psi_C$; no homogéneas: $\bK\bb\cdot\bn$ para $\psi_h$, inducida por el acoplamiento con el transporte. La condición terminal es nula porque $J$ es un
funcional integral sin término de costo terminal $\Phi(C(T))$; con un t\'ermino terminal aparecer\'ia $\psi_C(T)=\partial_C\Phi$.

%===========================================================
\section{Correspondencia con la implementaci\'on num\'erica}
\label{sec:codigo}
%===========================================================

\paragraph{Variantes.} Seis: \{ADR, MRMT semi-anal\'itico, MRMT por bloque\} $\times$ \{BFR (RBF-FD, \texttt{funtions/}), FEniCSx (\texttt{fenicsx/})\}. Par\'ametros: \texttt{factor\_sumidero\_C}$=0$ (no hay t\'ermino
expl\'icito de extracci\'on) y \texttt{factor\_almacenamiento\_C}$=1$ (t\'ermino $-S_s\partial_th\,C$).

\paragraph{Discretización del adjunto.} En BFR el adjunto es el del problema discreto (discretize-then-optimize); en FEniCSx se obtiene derivando los residuos con \texttt{ufl.derivative}. Notaci\'on discreta: $\mathcal{L}_h=J_h+\sum(\boldsymbol\psi^{n+1})^\top(\text{residuo})$, $J_h=\sum_{n=1}^{N_t}\gamma\Delta t(\mathbf{C}^n)^\top G\mathbf{C}^n$,
$\nabla j_n=2\gamma\Delta t\,G\mathbf{C}^n$. Los operadores $\mathcal{P},\mathcal{R}$ (y $\mathcal{P}_A,\mathcal{R}_A$) son los del Cap\'itulo~2 (trapezoidal, $\theta_C$), con la matriz diagonal de almacenamiento
\[
\Sigma^n=\frac{S_s}{\Delta t}\,\mathrm{diag}(\mathbf{h}^{n+1}-\mathbf{h}^n),
\]
que entra como $-\Sigma^n$ en la reacci\'on. Al ser diagonal es sim\'etrica.

\subsection*{ADR ($N_r=0$)}
\begin{equation}
(\mathcal{P}^{n})^\top\boldsymbol\psi_C^{n}=(\mathcal{R}^{n})^\top\boldsymbol\psi_C^{n+1}-\nabla j_n(\mathbf{C}^n),\quad n=N_t,\dots,1,\qquad\boldsymbol\psi_C^{N_t+1}=\mathbf{0}.
\label{eq:disc_adr}
\end{equation}

\subsection*{MRMT por bloque (complemento de Schur)}
$\mathcal{P}^{n+1}\mathbf{C}^{n+1}=\mathcal{R}^{n}\mathbf{C}^n+\sum_r\mathcal{Q}_r\mathbf{C}_{\im,r}^n+\mathbf{F}^n$ y $\mathbf{C}_{\im,r}^{n+1}=\kappa_r(\mathbf{C}^{n+1}+\mathbf{C}^n)+(1-2\kappa_r)\mathbf{C}_{\im,r}^n$
($\kappa_r=\Delta t\,\omega_r/\eta_r$, $\mathcal{Q}_r=4\phi R\,\omega_r I$). El adjunto:
\begin{equation}
\left\{\begin{aligned}
\boldsymbol\psi_r^{n}&=(1-2\kappa_r)\,\boldsymbol\psi_r^{n+1}+\mathcal{Q}_r\boldsymbol\psi_C^{n+1},\qquad \boldsymbol\psi_r^{N_t}=\mathbf{0},\\
(\mathcal{P}^{n})^\top\boldsymbol\psi_C^{n}&=(\mathcal{R}^{n})^\top\boldsymbol\psi_C^{n+1}+\sum_r\kappa_r\bigl(\boldsymbol\psi_r^{n}+\boldsymbol\psi_r^{n+1}\bigr)-\nabla j_n(\mathbf{C}^n).
\end{aligned}\right.
\label{eq:disc_block}
\end{equation}

\subsection*{MRMT semi-anal\'itico}
Con $e_r=e^{-\alpha_r\Delta t}$, $a_r=\eta_r(1-e_r)$: $\mathcal{P}_A^{n+1}\mathbf{C}^*=\mathcal{R}_A^{n}\mathbf{C}^n+\mathbf{F}^n$, $\tilde{\mathbf{C}}_r=(1-a_r)\tilde{\mathbf{C}}_{r-1}+a_r\mathbf{C}^n_{\im,r}$, $\mathbf{C}^{n+1}_{\im,r}=(1-e_r)\tilde{\mathbf{C}}_{r-1}+e_r\mathbf{C}^n_{\im,r}$, $\mathbf{C}^{n+1}=\tilde{\mathbf{C}}_{N_r}$, $\tilde{\mathbf C}_0=\mathbf C^*$.
\begin{equation}
\left\{\begin{aligned}
\boldsymbol\mu_{N_r}^{n+1}&=(\mathcal{R}_A^{n+1})^\top\boldsymbol\psi_C^{n+2}-\nabla j_{n+1}(\mathbf{C}^{n+1}),\\
\boldsymbol\mu_{r-1}^{n+1}&=(1-a_r)\,\boldsymbol\mu_r^{n+1}+(1-e_r)\,\boldsymbol\psi_r^{n+1},\quad r=N_r,\dots,1,\\
(\mathcal{P}_A^{n+1})^\top\boldsymbol\psi_C^{n+1}&=\boldsymbol\mu_0^{n+1},\qquad
\boldsymbol\psi_r^{n}=e_r\,\boldsymbol\psi_r^{n+1}+a_r\,\boldsymbol\mu_r^{n+1},
\end{aligned}\right.
\label{eq:disc_semi}
\end{equation}
con $\boldsymbol\psi_C^{N_t+1}=\boldsymbol\psi_r^{N_t}=\mathbf{0}$; para $N_r=0$ se recupera \eqref{eq:disc_adr}.

\subsection*{Adjunto discreto del flujo y gradiente}
Paso de flujo $A_h\mathbf{h}^{n+1}=B_h\mathbf{h}^n+\mathbf{f}^{n+1}$, $\mathbf{f}=-Q\mathbf{g}$. La fuente del adjunto del flujo tiene dos partes:
\begin{multline}
\boldsymbol\varpi^{m}=\underbrace{-\Bigl(\partial_{\mathbf{h}}\bigl[\mathcal{L}(\bu(\mathbf{h}^m))\mathbf{C}^{*,m}\bigr]\Bigr)^{\!\top}\boldsymbol\psi_C^{m}
-\Bigl(\partial_{\mathbf{h}}\bigl[\mathcal{L}(\bu(\mathbf{h}^m))\mathbf{C}^{m}\bigr]\Bigr)^{\!\top}\boldsymbol\psi_C^{m+1}}_{\text{velocidad}}\\
+\frac{S_s}{\Delta t}\Bigl[\boldsymbol\psi_C^{m+1}\!\circ(\mathbf{C}^{m+1}+\mathbf{C}^m)-\boldsymbol\psi_C^{m}\!\circ(\mathbf{C}^{m}+\mathbf{C}^{m-1})\Bigr]
\end{multline}
($\mathbf{C}^{*,m}=\mathbf{C}^m$ en ADR y bloque). La segunda parte es la derivada temporal discreta de $C\pA$, es decir, la discretizaci\'on de $-S_s\partial_t(C\pA)$ de \eqref{eq:adj_flujo_A}. Entonces
\begin{equation}
A_h^\top\boldsymbol\psi_h^{m}=B_h^\top\boldsymbol\psi_h^{m+1}-\boldsymbol\varpi^{m},\quad\boldsymbol\psi_h^{N_t+1}=\mathbf{0},\qquad
\frac{\partial\hat J}{\partial Q_n}=2\beta d_p\Delta t\,Q_n+\mathbf{g}^\top\boldsymbol\psi_h^{n+1}\ (\text{con pesos }\theta).
\label{eq:disc_flujo}
\end{equation}

\paragraph{Estado de la implementaci\'on.} Las ecuaciones de estado y el adjunto $\boldsymbol\psi_C$ est\'an implementados en BFR (\texttt{funtions/operators.py}, \texttt{upwind\_operators.py}, \texttt{step\_time.py}) y FEniCSx (\texttt{fenicsx/forward.py},
\texttt{adjoint.py}) para las tres variantes de modelo. Las discrepancias detectadas entre formulaci\'on y c\'odigo, y las pruebas efectivamente ejecutadas, se resumen en la Secci\'on~\ref{sec:verif} y en el informe de auditor\'ia (\texttt{docs/auditoria\_reanalisis.md}).

%===========================================================
\section{Verificaciones matem\'aticas}
\label{sec:verif}
%===========================================================

\paragraph{Qu\'e se verific\'o y c\'omo.} Las comprobaciones son de dos tipos: las \emph{anal\'iticas} se
demuestran en el texto (se indica la secci\'on) y las \emph{num\'ericas} se ejecutaron con
\texttt{tests/test\_gradiente\_modelo\_nuevo.py}. Lo que no se ejecut\'o se declara.

\begin{center}\small
\begin{tabular}{p{3.6cm}p{6.6cm}p{4.2cm}}
\toprule
Verificaci\'on & Estado & D\'onde\\
\midrule
Dimensional & anal\'itica: todos los t\'erminos de \eqref{eq:N_ADR} y \eqref{eq:N_MRMT} tienen unidades $[C]/$s; unidades de $\psi$ y gradientes en la tabla de notaci\'on & \S\ref{sec:gobierno}, tabla de notaci\'on\\
Signos (difusi\'on, advecci\'on, reacci\'on, intercambio) & anal\'itica: transposici\'on t\'ermino a t\'ermino & \S\ref{sec:adj_ADR}, \S\ref{sec:adj_MRMT}\\
Conservaci\'on de masa & anal\'itica: \eqref{eq:balance} y MRMT (el intercambio se cancela). \textbf{No se ejecut\'o num\'ericamente en esta sesi\'on} & \S\ref{sec:gobierno}, \S\ref{sec:mrmt}\\
Frontera & anal\'itica: cada t\'ermino de borde de las formas d\'ebiles se anula o se conserva por una condici\'on de la tabla de fronteras & \S\ref{sec:debil}, \S\ref{sec:bc_adj}\\
Variacional (adjunto = derivada del Lagrangiano) & anal\'itica y num\'erica (tablas siguientes) & \S\ref{sec:adj_ADR}--\ref{sec:adj_MRMT}\\
Terminal & anal\'itica: $J$ sin t\'ermino terminal implica $\psi(T)=0$ & \S\ref{sec:bc_adj}\\
Equivalencia conservativa/expandida & anal\'itica: Proposici\'on~\ref{prop:equiv} & \S\ref{sec:gobierno}\\
\bottomrule
\end{tabular}
\end{center}

\paragraph{Gradiente adjunto frente a diferencias finitas centradas: CI uniforme $C\equiv1$, con almacenamiento.}
Error relativo m\'aximo sobre cuatro direcciones ($-\nabla J$, $Q$ uniforme, $Q$ aleatoria, $(Q,z_p)$ aleatoria) y dos pasos $\varepsilon$; s\'olo se compara $\gamma J_m$ (la parte econ\'omica es anal\'itica y se resta).
\begin{center}\small\begin{tabular}{lcc}\toprule Modelo & BFR & FEniCSx\\\midrule
ADR & no ejecutado & $8.6\times10^{-7}$\\
MRMT semi-anal\'itico & no ejecutado & $5.8\times10^{-6}$\\
MRMT por bloque & no ejecutado & $9.2\times10^{-7}$\\
\bottomrule\end{tabular}\end{center}
\paragraph{Gradiente adjunto frente a diferencias finitas centradas: CI modulada $C=1+\tfrac12\cos(2\pi x/2000)\cos(2\pi z/800)$, con almacenamiento.}
Error relativo m\'aximo sobre cuatro direcciones ($-\nabla J$, $Q$ uniforme, $Q$ aleatoria, $(Q,z_p)$ aleatoria) y dos pasos $\varepsilon$; s\'olo se compara $\gamma J_m$ (la parte econ\'omica es anal\'itica y se resta).
\begin{center}\small\begin{tabular}{lcc}\toprule Modelo & BFR & FEniCSx\\\midrule
ADR & $5.9\times10^{0}$ & $9.0\times10^{-7}$\\
MRMT semi-anal\'itico & no ejecutado & $5.4\times10^{-6}$\\
MRMT por bloque & no ejecutado & $9.5\times10^{-7}$\\
\bottomrule\end{tabular}\end{center}
\paragraph{Gradiente adjunto frente a diferencias finitas centradas: CI modulada, sin almacenamiento.}
Error relativo m\'aximo sobre cuatro direcciones ($-\nabla J$, $Q$ uniforme, $Q$ aleatoria, $(Q,z_p)$ aleatoria) y dos pasos $\varepsilon$; s\'olo se compara $\gamma J_m$ (la parte econ\'omica es anal\'itica y se resta).
\begin{center}\small\begin{tabular}{lcc}\toprule Modelo & BFR & FEniCSx\\\midrule
ADR & $1.0\times10^{0}$ & $1.2\times10^{-6}$\\
MRMT semi-anal\'itico & no ejecutado & $2.8\times10^{-6}$\\
MRMT por bloque & no ejecutado & $9.8\times10^{-7}$\\
\bottomrule\end{tabular}\end{center}

Condiciones: dominio real, 3 pasos de $\Delta t=36000$\,s, $\theta_C$ del entorno de prueba, pozo y caudal por defecto.
El adjunto de FEniCSx se obtiene derivando los residuos discretos (\texttt{ufl.derivative}, salvo las recursiones MRMT escritas a mano) y
debe coincidir con las diferencias finitas hasta el redondeo; el de BFR es la transpuesta de los operadores discretos y
est\'a sujeto, adem\'as, al tratamiento del nodo del pozo y de los nodos de frontera.


\paragraph{Limitaciones del adjunto BFR.} Con concentraci\'on no uniforme el adjunto de BFR no coincide con las diferencias finitas
(factor 2--7 con almacenamiento; sin almacenamiento el adjunto es $\approx0$ y las diferencias finitas dan valores grandes). No cambia
al desactivar el recorte $C\ge0$, al imponer la transpuesta exacta en $\Gin\cup\Gout$ ni al a\~nadir el t\'ermino $\partial_{\bu}\bD$.
Se atribuye, sin haberlo demostrado, a una sensibilidad espuria del forward en el nodo del pozo; \textbf{el adjunto BFR no est\'a validado}.

\paragraph{Comparaci\'on BFR--FEniCSx del modelo directo.} Dominio real, 30 d, $\Delta t=18000$\,s, mismo mallado y misma CI, sin $W_-$ y con
almacenamiento. Error relativo $L^2$ final de $C$ entre BFR y FEniCSx: ADR $0.27$, MRMT semi-anal\'itico $0.27$, MRMT por bloque $0.27$;
cociente de masas BFR/FEniCSx $0.94$ en las tres. La diferencia procede de las discretizaciones y de la condici\'on de salida de flujo (Cap.~4 de la tesis); no se
atribuy\'o aqu\'i a ning\'un t\'ermino concreto.


%===========================================================
\section{Tabla consolidada de notaci\'on}
%===========================================================

La tabla de la Secci\'on~3 es la tabla consolidada. A los s\'imbolos all\'i listados se a\~naden los discretos: $\mathcal{P},\mathcal{R}$ (operadores impl\'icito y expl\'icito), $\Sigma^n$ (almacenamiento),
$\kappa_r,\mathcal{Q}_r$ (acoplamiento por bloque), $e_r,a_r$ (semi-anal\'itico), $\boldsymbol\psi_r^n,\boldsymbol\mu_r^n$ (adjuntos inm\'ovil e intermedio), $\boldsymbol\varpi^m$ (fuente discreta del adjunto del flujo), $A_h,B_h$ (matrices del flujo).

%===========================================================
\section{Conclusiones}
%===========================================================

La extracci\'on del pozo queda incorporada de forma impl\'icita en el transporte a trav\'es de $\nabla\!\cdot\bu$; el balance de masa \eqref{eq:balance} la recupera como $-Q\int\geps C$. El transporte conserva
$-S_s\partial_th\,C$, que se cancela en el adjunto en forma no conservativa pero reaparece como fuente $-S_s\partial_t(C\psi)$ en el adjunto del flujo. Los adjuntos de ADR y MRMT, con sus condiciones de frontera naturales
$(\bD\nabla\psi+\bu\psi)\cdot\bn=0$, y el gradiente respecto a $Q$ y $z_p$ se obtienen por variaci\'on de los funcionales de Lagrange. El adjunto continuo del flujo est\'a sobredeterminado en $\Gout$ y por eso se usa el adjunto discreto.

%===========================================================
\begin{thebibliography}{9}
%===========================================================
\bibitem{Evans2010} L.~C.~Evans, \emph{Partial Differential Equations}, 2.\textordmasculine\ ed., AMS, 2010.
\bibitem{LionsMagenes1972} J.-L.~Lions y E.~Magenes, \emph{Non-Homogeneous Boundary Value Problems and Applications}, vol.~I, Springer, 1972.
\bibitem{DautrayLions1992Vol5} R.~Dautray y J.-L.~Lions, \emph{Mathematical Analysis and Numerical Methods for Science and Technology}, vol.~5, Springer, 1992.
\bibitem{Troltzsch2010} F.~Tr\"oltzsch, \emph{Optimal Control of Partial Differential Equations}, AMS, 2010.
\bibitem{Hinze2009} M.~Hinze, R.~Pinnau, M.~Ulbrich y S.~Ulbrich, \emph{Optimization with PDE Constraints}, Springer, 2009.
\bibitem{bear1972dynamics} J.~Bear, \emph{Dynamics of Fluids in Porous Media}, Elsevier, 1972.
\bibitem{Haggerty1995} R.~Haggerty y S.~M.~Gorelick, ``Multiple-rate mass transfer for modeling diffusion and surface reactions in media with pore-scale heterogeneity'', \emph{Water Resour.\ Res.} 31(10), 1995.
\end{thebibliography}

\end{document}
